\documentclass{article}
\usepackage[utf8]{inputenc}
\usepackage{hyperref}
\usepackage{fancyvrb}
\usepackage[a4paper, total={6in, 8in}]{geometry}

\geometry{legalpaper, portrait, margin=1in}

\hypersetup{
    colorlinks = true,
    allcolors = black,
    urlcolor = blue,
}

\setlength{\parindent}{20pt}
\title{INE5416 - Paradigmas de Programação (2016/2) - Trabalho 3}
\author{
	Bruno Marques do Nascimento\\
	\texttt{brunomn95@gmail.com}\\
	\\ 
	Gustavo Figueira Olegário\\
	\texttt{gustavo-olegario@hotmail.com}\\
	\\
	João Victor de Mello Fagundes\\
	\texttt{joaovictordmf@gmail.com}\\
	\\
}
\date{\today}

\usepackage{natbib}
\usepackage{graphicx}
\usepackage{caption}
\begin{document}
\maketitle

\noindent\rule{16cm}{0.4pt}

\section*{Imagens fornecidas no arquivo 'imgs.zip':}
\paragraph{}

\begin{center}
\includegraphics[scale=1.00]{texImgs/img1.eps}
\includegraphics[scale=1.00]{texImgs/img2.eps}
\includegraphics[scale=1.00]{texImgs/img3.eps}
\includegraphics[scale=1.00]{texImgs/img4.eps}
\includegraphics[scale=1.00]{texImgs/img5.eps}
\end{center}

\begin{center}
\includegraphics[scale=1.00]{texImgs/img6.eps}
\includegraphics[scale=1.00]{texImgs/img7.eps}
\includegraphics[scale=1.00]{texImgs/img8.eps}
\includegraphics[scale=1.00]{texImgs/img9.eps}
\end{center}

\section*{Implementações solicitadas:}
\subsection*{Calculador de centroide}
\paragraph{}

Para utilizá-lo basta utilizar o comando:

\begin{center}
\begin{minipage}{13cm}
\begin{Verbatim}

?- centroid('imgs/img1.pgm', (X,Y)).
X = 18,
Y = 19.
\end{Verbatim}
\end{minipage}
\end{center}

É possível gerar uma imagem com seu respectivo centroide destacado, com o seguinte comando:
	
\begin{center}
\begin{minipage}{13cm}
\begin{Verbatim}

?- centroid('imgs/img1.pgm', _, CentroidFileName, CentroidMatrix), 
   writePGM(CentroidFileName, FinalCentroidM). 
\end{Verbatim}
\end{minipage}
\end{center}

\begin{center}
\includegraphics[scale=1.00]{texImgs/centroid_img1.eps}
\end{center}
	

Além disso, existe um comando para gerar todas as imagens fornecidas pelo professor no arquivo imgs.zip, com seus respectivos centroids.

\begin{center}
\begin{minipage}{13cm}
\begin{Verbatim}

?- generateAllCentroid().

\end{Verbatim}
\end{minipage}
\end{center}

\begin{center}
\includegraphics[scale=1.00]{texImgs/centroid_img2.eps}
\includegraphics[scale=1.00]{texImgs/centroid_img3.eps}
\includegraphics[scale=1.00]{texImgs/centroid_img4.eps}
\includegraphics[scale=1.00]{texImgs/centroid_img5.eps}
\end{center}

\begin{center}
\includegraphics[scale=1.00]{texImgs/centroid_img6.eps}
\includegraphics[scale=1.00]{texImgs/centroid_img7.eps}
\includegraphics[scale=1.00]{texImgs/centroid_img8.eps}
\includegraphics[scale=1.00]{texImgs/centroid_img9.eps}
\end{center}


\subsection*{Detector de borda}
\paragraph{}
	Para utilizar o detector de bordas, ou seja, gerar a imagem que representa a borda da imagem fornecida, basta utilizar o seguinte comando:

\begin{center}
\begin{minipage}{13cm}
\begin{Verbatim}

?- border('imgs/img1.pgm', _, _).
\end{Verbatim}
\end{minipage}
\end{center}

\begin{center}
\includegraphics[scale=1.00]{texImgs/border_img1.eps}
\end{center}

	Também é possível gerar a borda de todas as imagens automaticamente, através do comando:

\begin{center}
\begin{minipage}{13cm}
\begin{Verbatim}

?- generateAllBorder().

\end{Verbatim}
\end{minipage}
\end{center}	

\begin{center}
\includegraphics[scale=1.00]{texImgs/border_img2.eps}
\includegraphics[scale=1.00]{texImgs/border_img3.eps}
\includegraphics[scale=1.00]{texImgs/border_img4.eps}
\includegraphics[scale=1.00]{texImgs/border_img5.eps}
\end{center}

\begin{center}
\includegraphics[scale=1.00]{texImgs/border_img6.eps}
\includegraphics[scale=1.00]{texImgs/border_img7.eps}
\includegraphics[scale=1.00]{texImgs/border_img8.eps}
\includegraphics[scale=1.00]{texImgs/border_img9.eps}
\end{center}


\subsection*{Extrator de característica}
\paragraph{}
	O extrator de característica fornece a distância média entre os pontos de uma determinada imagem e seu centróide, e o desvio padrão destas distâncias. Ele é utilizado com a seguinte chamada:

\begin{center}
\begin{minipage}{13cm}
\begin{Verbatim}

?- characteristicExtractor('imgs/img1.pgm', MeanDistance, StandardDeviation).
   MeanDistance = 9.502553256112856,
   StandardDeviation = 3.379121593967018.

\end{Verbatim}
\end{minipage}
\end{center}	

\subsection*{Reconhecedor de imagem com círculo}
\paragraph{}
	Para este reconhecedor de imagem com círculo, partiu-se do prícipio que o desvio padrão da distância dos pontos presentes na borda de uma circunferência ao seu centro sempre será o menor dentre as formas conhecidas(Ex: Quadrado, Triangulo, etc.), pois é na circunferência onde temos a menor variação de distância já que todos estão a uma mesma distância do centro. Com isso foi definido um valor inicial de desvio padrão de 0.54 no arquivo 'database.pl', que será alterado conforme o usuário identifica círculo nas imagens fornecidas. Abaixo temos um exemplo de uso do reconhecedor de imagem:
	\\
\begin{center}
\includegraphics[scale=1.00]{texImgs/img7.eps}
\end{center}

\begin{center}
\begin{minipage}{6cm}
\begin{Verbatim}
?- isCircle('imgs/img7.pgm').
Identifying...This is a circle!


\end{Verbatim}
\end{minipage}
\end{center}

\begin{center}
\includegraphics[scale=1.00]{texImgs/img8.eps}
\end{center}

\begin{center}
\begin{minipage}{6cm}
\begin{Verbatim}
?- isCircle('imgs/img8.pgm').
Identifying...This is not a circle!
Is this a circle? [y/n]y.


\end{Verbatim}
\end{minipage}
\end{center}

\begin{center}
\includegraphics[scale=1.00]{texImgs/img8.eps}
\end{center}

\begin{center}
\begin{minipage}{6cm}
\begin{Verbatim}
?- isCircle('imgs/img8.pgm').
Identifying...This is a circle!
\end{Verbatim}
\end{minipage}
\end{center}

\end{document}